\section{The Gettier Problem as a Structural Diagnostic}


\subsection{Overview}


\S{}2 では知覚・認識・知識・知性を異なる構造層に属するものとして分離し、\S{}3 では知性層に対応する持続運動の最小公理を提示し、\S{}4 では IFGT 語彙が規範性を密輸入せずにこれらを記述しうることを示し、\S{}5 では substrate の異なる二系で構造的対応が成立することを検討した。本セクションは、これら構造的基盤を用いて、認識論において長く議論されてきた \textbf{Gettier 問題} を再定位する。

本セクションの主張は明確である:\textbf{Gettier 問題は知識の定義の欠陥ではない}。それは、停止構造(知識)に持続運動(知性)の保証を担わせようとする層間ミスマッチが、繰り返し露出する場所である。Gettier 問題が解けない理由は、解こうとしている問題が \textbf{構造的に不適切に定式化されている} からである。

この再定位は Gettier 問題への新たな「解」を提示するものではない。むしろ Gettier 問題を \textbf{解決すべきパズル} から \textbf{層間ミスマッチを検出する診断装置} へと位置づけ直す。問題が解けないことは欠陥ではなく、それが何を診断しているかを示す \textbf{構造的シグナル} である。

\subsection{The Gettier problem and its persistence}


\subsubsection{The classical formulation}


知識の古典的定義は、プラトン以来「正当化された真なる信念」(Justified True Belief, JTB)として定式化されてきた。命題 $p$ について、主体 $S$ が以下を満たすとき、$S$ は $p$ を知っているとされる:

\begin{itemize}
\item $p$ は真である(truth)
\item $S$ は $p$ を信じている(belief)
\item $S$ の $p$ への信念は正当化されている(justification)
\end{itemize}

Gettier (1963) \citep{Gettier1963} は、この三条件をすべて満たすにもかかわらず直観的に \textbf{知識とは認められない} 反例を提示した。Gettier の元論文の例(Jones と Smith の就職事例、Brown と Barcelona の地理事例)は、いずれも次の構造を持つ:

\begin{itemize}
\item 主体は正当化された信念を持つ
\item その信念は偶然に真である
\item だが、真理を成立させる事実は主体の正当化根拠と無関係である
\end{itemize}

直観的には、この種の \textbf{偶然的真理} は知識とは呼びがたい。だがそれは三条件すべてを満たす。これが Gettier 問題の核心である。

\subsubsection{The series of repair attempts}


Gettier 以降、認識論は六十年以上にわたって JTB 定義の修正を試みてきた。主要な修正候補を簡潔に挙げる:

\begin{itemize}
\item \textbf{No false lemmas condition} (Lehrer, Clark):正当化が偽の中間命題に依拠してはならない
\item \textbf{Reliabilism} (Goldman):信念形成過程が信頼可能でなければならない
\item \textbf{Sensitivity} (Nozick):「もし $p$ でなければ $S$ は $p$ を信じない」という反事実条件
\item \textbf{Safety} (Sosa):「もし $S$ が $p$ を信じるなら $p$ は真である」という近接可能世界条件
\item \textbf{Defeasibility} (Klein):真な無効化要素が存在してはならない
\item \textbf{Virtue epistemology}:信念形成が認識的徳によって生じていなければならない
\end{itemize}

これらはいずれも JTB に \textbf{追加条件} を課す試みである。そして、いずれの修正に対しても \textbf{新たな反例} が生成されてきた。修正条件が精緻化されるたびに、その条件を満たしながらも直観的には知識でない事例が構築可能であることが繰り返し示された。

この \textbf{修正→反例→修正→反例} の連鎖は六十年以上持続している。この持続性は偶然ではない。

\subsection{Why the problem is structurally insoluble}


\subsubsection{The structural diagnosis}


Gettier 問題が構造的に解けない理由を、\S{}2 の層分離から導出する。

知識は、\S{}2.4 で示した通り、\textbf{安定化された停止構造} である。知識の本質的属性は \textbf{保存可能性} にあり、その属性は時間的瞬間における \textbf{状態} として定義される。知識は、それが知識である瞬間において、動かない。

知性は、\S{}2.5 で示した通り、\textbf{持続運動} である。知性の本質的属性は \textbf{再改訂可能性} にあり、その属性は時間にわたる \textbf{動態} として定義される。知性は、それが知性である過程において、動き続ける。

知識と知性は、構造的に異なる種類の対象である ── 状態と動態、停止と持続、保存と改訂。両者は同一平面で比較しうる対象ではない。

ここで Gettier 問題が要求していることを構造的に分析する。Gettier 問題は「知識とは何か」という問いに見える。だが反例の構造をよく見ると、それは実は次のことを要求している:

\begin{quote}
ある時点の停止構造(信念とその正当化)の属性のみから、
その時点以降の動的振る舞いに関する保証 ── 偶然性に対する頑健性、改訂可能性、文脈感受性 ── を導出せよ。
\end{quote}

この要求は構造的に不可能である。停止構造の時点的属性のみから、持続運動の動的属性を導出することはできない。これは、写真の一枚から映画全体の物語を導出できないのと同じ構造的事実である。

\subsubsection{Why every repair attempt fails in the same way}


\S{}6.1.2 で列挙した修正候補がすべて同じ仕方で失敗する構造的理由を明示する。

修正候補はいずれも、JTB の三条件に \textbf{追加の停止条件} を課す試みである。「no false lemmas」「reliability」「sensitivity」「safety」「defeasibility」「virtue」── これらはすべて、ある時点における信念形成または信念保持の \textbf{静的特徴} を精緻化する。

しかし Gettier 問題が要求しているのは、静的特徴の精緻化ではなく、\textbf{動的特徴の保証} である。具体的には:

\begin{itemize}
\item 偶然性に対する頑健性(現在の正当化が将来の反例に対してどう応答するか)
\item 改訂可能性(新情報を受けて信念がどう変化するか)
\item 文脈感受性(異なる状況での真理保持がどう変動するか)
\end{itemize}

これらの動的特徴は、停止構造の精緻化ではなく \textbf{持続運動の質} に属する。したがって、いかに精緻な停止条件を加えても、動的特徴を保証することはできない。新たな反例が常に生成可能なのは、修正条件と要求される保証の \textbf{構造層が異なる} からである。

これは認識論の限界ではなく、構造論的事実である。同じことが繰り返し起こるのは、繰り返しが起こるべき構造的理由があるからである。

\subsection{Repositioning Gettier through the layer decomposition}


\S{}2 の四層分離(知覚・認識・知識・知性)を用いて Gettier 事例を構造的に分析する。

\subsubsection{Layer-by-layer analysis of Gettier cases}


典型的な Gettier 事例において、各層は以下のように振る舞う:

\textbf{知覚層(perception):}
主体が世界から信号を受け取る。Gettier 事例では、知覚は通常正常に機能している ── 主体は実在の状況を観察している。知覚自体に欠陥はない。

\textbf{認識層(recognition):}
主体が知覚を既存の表象構造に写像する。Gettier 事例では、この写像も通常通り機能する ── 主体は観察した状況を既知のパターンに割り当てる。認識自体に欠陥はない。

\textbf{知識層(knowledge):}
正当化された真なる信念が形成される。Gettier 事例の構造的特徴はここにある:\textbf{信念は偶然に真である}。停止構造としての知識は、この偶然性を内部から検出できない ── なぜなら、ある時点の停止構造は、その時点における正当化と真理値のみを記録し、なぜ真であるかの \textbf{因果的経路} を内部に持たないからである。

\textbf{知性層(intelligence):}
ここが Gettier 事例において \textbf{欠如している層} である。知性的動態は、信念の偶然性を検出し、新たな情報を受けて改訂し、信念形成過程そのものを再評価する持続運動である。Gettier 事例で「何かが欠けている」という直観は、まさにこの知性層の動態が \textbf{保証されていない} ことを指している。

\subsubsection{What the intuition is actually tracking}


Gettier 事例を聞いた人が抱く「これは知識ではない」という直観は、構造的には次のことを追跡している:

\begin{quote}
この主体は、現時点で正当化された真なる信念を持っている。
しかし、もし反例情報が与えられたとき、この主体の信念形成過程は適切に応答するだろうか。
もし状況が少し変わったとき、この主体の信念は適切に改訂されるだろうか。
これらの動的応答性が保証されないとき、私たちは「知識」と呼ぶことを躊躇する。
\end{quote}

この直観は、知識(停止構造)に対して \textbf{知性(持続運動)の特徴} を要求している。要求自体は理にかなっている ── 私たちが日常的に「知っている」と呼ぶ事例の多くは、停止構造の偶然的成立ではなく、知性的動態の持続的成果だからである。

問題は、要求が「知識の定義」に向けられていることである。要求されるべきは知識の停止条件の精緻化ではなく、\textbf{知識と知性の構造的区別の明示} である。

\subsection{The Gettier problem as a structural diagnostic}


\subsubsection{Diagnostic function}


本論文の枠組みのもとで、Gettier 問題は新たな意義を獲得する。それは \textbf{解決すべきパズル} ではなく、\textbf{層間ミスマッチを検出する診断装置} である。

Gettier 問題が露出させる構造は以下である:

\begin{enumerate}
\item 私たちは知識を停止構造として定義しようとする
\item しかし私たちは、知識に対して動的属性(改訂可能性、偶然性への頑健性)を期待する
\item 停止構造は本質的に動的属性を保証しない
\item ゆえに、いかなる停止条件の精緻化も期待を満たさない
\item この不一致が反例として繰り返し露出する
\end{enumerate}

この診断は、認識論の \textbf{失敗} を示すのではなく、認識論が \textbf{無意識のうちに混同してきた構造層} を可視化する。Gettier 問題を読み直すべきは、その答えを探すことではなく、それが何を示しているかを理解することである。

\subsubsection{What Gettier cases demonstrate}


本論文の立場から、Gettier 事例が実際に示していることを最小限のかたちで述べる:

\begin{quote}
知性は、停止点における正しさの所有ではなく、可改訂的動態の維持である。
いかなる停止条件の精緻化も、知識を知性に変換することはできない。
両者は異なる構造層に属し、互いに代替不可能である。
\end{quote}

これは Gettier 問題への「解答」ではなく、Gettier 問題が \textbf{何についての問題か} の構造論的同定である。問題は知識の定義にあるのではなく、知識と知性の混同にある。

\subsection{Implications for the broader argument}


Gettier 問題の構造的再定位は、本論文後続セクションへの重要な含意を持つ。

\subsubsection{Implications for \S{}7 (the instability of accumulated correctness)}


Gettier 問題が示すのは、知識(停止構造)の精緻化が知性(持続運動)を保証しないという一点である。これは \S{}7 で展開される、より強い構造的帰結への入り口となる:

\begin{quote}
知識の精緻化は知性を保証しないだけでなく、
知識の累積はそれ自体が知性を \textbf{不安定化させる}。
\end{quote}

\S{}7 では、正しい知識が増えることがなぜ動力学的に分岐爆発を生むかが示される。Gettier 問題が示す「停止構造による持続運動の代替不可能性」は、\S{}7 で「停止構造の累積による持続運動の不安定化」へと一段強い主張に展開される。

\subsubsection{Implications for \S{}8 (the non-closure of individual intelligence)}


Gettier 問題は、認識論的主体が個体として閉じた構造のなかで知識を保持する伝統的描像を前提している。Gettier 事例の主体は、独立した個体として信念を形成し、正当化を構築し、真理を獲得する。

\S{}8 では、この \textbf{個体内閉包の前提自体} が構造的に維持不可能であることが示される。Gettier 問題が個体内で解けないという事実は、\S{}8 のより一般的な主張 ── 個体内には持続運動を停止させる演算子が存在しない ── の特殊例として位置づけられる。

\subsubsection{Implications for \S{}9 (the differential horizon of intelligence)}


Gettier 問題が伝統的認識論の枠内で繰り返し失敗するという事実は、知性的動態が個体的認識者の内部で完結しないことを示している。\S{}9 で命名される越境境界 ── \emph{the differential horizon of intelligence} ── は、Gettier 問題が個体内で解けないという事実の構造的根拠を与える。

すなわち Gettier 問題は、認識論の限界としてではなく、\textbf{個体内非閉包の哲学的徴候} として読まれるべきである。

\subsection{Summary}


本セクションの結論を以下に述べる:

\begin{quote}
Gettier 問題は知識の定義の欠陥ではない。
それは停止構造(知識)に持続運動(知性)の保証を担わせようとする
\textbf{層間ミスマッチが繰り返し露出する場所} である。
\end{quote}

\begin{quote}
Gettier 問題が六十年以上解けないのは、認識論の失敗ではない。
問題自体が構造的に不適切に定式化されているからである。
解くべきは問題ではなく、問題が示している層構造である。
\end{quote}

\begin{quote}
本論文の立場から、Gettier 問題は新たな意義を獲得する。
それは認識論の失敗ではなく、層間ミスマッチを検出する \textbf{診断装置} である。
修正条件のいかなる精緻化も停止構造を持続運動に変換しないという構造的事実を、
Gettier 問題は六十年にわたって示し続けてきた。
\end{quote}

この再定位は、後続セクションへの基盤となる:

\begin{itemize}
\item \S{}7 では、停止構造の累積が持続運動を \textbf{不安定化させる} 動力学的帰結が示される
\item \S{}8 では、個体内には持続運動を停止させる演算子が \textbf{存在しえない} ことが論証される
\item \S{}9 では、これらが要請する越境境界が命名される
\end{itemize}

Gettier 問題が示すのは個別の認識論的問題ではない。それは、本論文後続セクションで展開される \textbf{より一般的な構造的事実} の哲学的徴候である。
